\documentclass{article}
\usepackage[utf8]{inputenc}
\usepackage[T1]{fontenc}
\usepackage{geometry}
\usepackage{graphicx}
\usepackage{float}
\usepackage{hyperref}
\usepackage[english,indonesian]{babel}

\title{Optimasi Pemrosesan Template LaTeX Menggunakan Pendekatan Berbasis AI}
\author{Academic Support Agent}
\date{\today}

\begin{document}
\maketitle

\begin{abstract}

\end{abstract}

\chapter{PENDAHULUAN}


\chapter{TINJAUAN PUSTAKA}


% Original content (commented out):
% %Esta versión en LaTeX de formato de informes de practica fue hecho por Mario Dinamarca a partir del modelo de informes en LaTeX que fue publicado por Maximiliano Muñoz Leiva en el grupo de Facebook de los estudiantes de la escuela de ingeniería civil industrial.

\documentclass[12pt,letter,spanish]{article}
\usepackage[left=3cm,right=2.5cm,top=3cm,bottom=2.5cm]{geometry}
\usepackage[utf8]{inputenc}
\usepackage[spanish]{babel}
\usepackage{times} %Para que salga en times
\usepackage[breakable]{tcolorbox}
\usepackage{dirtytalk}
\usepackage{seqsplit}
\usepackage{wrapfig}
\usepackage{enumitem}
\usepackage{float}
\usepackage[hidelinks]{hyperref}
\usepackage{fancyhdr}
\usepackage{comment}

\newcommand{\source}[1]{\caption*{Fuente: {#1}} }%Source a figuras



%%%%%%%%%%%%%%%%%%%%%%%%%%%%%%%%%%%% CONFIGURACIÓN ÍNDICE %%%%%%%%%%%%%%%%%%%%%%%%%%%%%%%%%%%%%%%%%


\usepackage[nottoc,numbib]{tocbibind}% make LoF and LoT show up in ToC
\usepackage[titles,subfigure]{tocloft} %%%Para tabla de contenidos
%%%%%%%% Alinear indice a la izquierda
\setlength{\cftpartindent}{0em}
\setlength{\cftsecindent}{0em}
\setlength{\cftsubsecindent}{0em}
\setlength{\cftsubsubsecindent}{0em}

%%%%%%%%%%%%%%%%%%%% CONFIGURACIÓN LISTA TABLAS Y FIGURAS %%%%%%%%%%%%%%%%%%%%
\usepackage[labelfont=bf]{caption}%Nombre figuras y tablas en negrita

\renewcommand{\cftsecleader}{\cftdotfill{\cftdotsep}}

\renewcommand{\cftfigfont}{Figura }%Lo que sale en los índices tablas
\renewcommand{\cfttabfont}{Tabla }%Lo que sale en los índices tablas

\renewcommand{\cftfigaftersnum}{\hfill--\hfill} 
\renewcommand{\cfttabaftersnum}{\hfill--\hfill}

\renewcommand{\cftfignumwidth}{3em}
\renewcommand{\cfttabnumwidth}{3em}



\usepackage[toc]{glossaries} %Agregar glosarios a index
\usepackage{caption}%Source a figuras

%%%%%%%%%%%%%%%%%%%%   REFERENCIAS  %%%%%%%%%%%%%%%%%%%%
\usepackage[backend=biber, style=apa]{biblatex}
\def\bibtitle{Referencias}
\addbibresource{bibliografia.bib}


\renewcommand{\headrulewidth}{0pt}
% your new footer definitions here


\usepackage{graphicx}
\usepackage{caption}
\usepackage{amsmath}
\usepackage{longtable}
\usepackage{subfigure} 
\usepackage{ragged2e}
\usepackage{multicol}
\usepackage{xcolor}
\usepackage{rotating}
\usepackage{xcolor}
\usepackage{sectsty}
\usepackage{glossaries}
\usepackage{tikz}
\newcommand{\q}[1]{``#1''} %Para texto entre comillas

\renewcommand{\baselinestretch}{1.15}

%%%%%%%%%%%%%%%%%%%%%%%%%%%%%%%%%%%%%%%% INICIO GLOSARIO %%%%%%%%%%%%%%%%%%%%%%%%%%%%%%%%%%%%%%%%%%

\makeglossaries

%%%%%%%%%%%%%%%%%%%%%%%%%%%%%%%%%%%%%%%% FIN GLOSARIO %%%%%%%%%%%%%%%%%%%%%%%%%%%%%%%%%%%%%%%%%%

\begin{document}

\input{portada}


\renewcommand{\listtablename}{Lista de Tablas}%%%%%%%%%%% cambio de nombre
\renewcommand{\listfigurename}{Lista de Figuras} %%%%%%% cambio de nombre
\renewcommand{\tablename}{Tabla}

\numberwithin{figure}{section}%Numeración multinivel en índice
\numberwithin{table}{section} %Numeración multinivel en índice

\fancyhf{}

\fancyhead[C]{\thepage}

\pagestyle{fancy}

\fancypagestyle{plain}

\fancyhf{} 

\fancyhead[C]{\thepage}



\setlength{\parindent}{1cm} %%%%%%%% Sangría al inicio de cada párrafo
\setlength{\parskip}{1em}

%%%%%%%%%%%%%%%%%%%%%%%%%%%%%%%%%%%%%%%% INICIO DOCUMENTO %%%%%%%%%%%%%%%%%%%%%%%%%%%%%%%%%%%%%%%%%%

%%%%%%%%%%%%%%%%%%%%%%%%%%%%%%%%%%%%%%%%  INICIO ÍNDICE   %%%%%%%%%%%%%%%%%%%%%%%%%%%%%%%%%%%%%%%%%%
\newpage

\begin{center}%%%%%% inicio centrado índice, figuras, tablas

\tableofcontents
\clearpage

%%%%%%%%%%%%%%%%%%%%%%%%%%%%%%%%%%%%%%%% INICIO GLOSARIO %%%%%%%%%%%%%%%%%%%%%%%%%%%%%%%%%%%%%%%%%%


\input{glosario}
\setcounter{tocdepth}{4}
\setcounter{secnumdepth}{4}

\glsaddall
\glsaddallunused
\printglossary[type=main, title=Glosario]

%%%%%%%%%%%%%%%%%%%%%%%%%%%%%%% INICIO LISTA FIGURAS Y TABLAS %%%%%%%%%%%%%%%%%%%%%%%%%%%%%%%%%%%%%

\clearpage
\listoffigures

\clearpage

\listoftables
\clearpage

\end{center}%%%%%%% fin centrado índice, figuras, tablas

%%%%%%%%%%%%%%%%%%%%%%%%%%%%%%%%%%%%%%%%%%%%%%%%%%%%%%%%%%%%%%%%%%%%%%%%%%%%%%%%%%%%%%%%%%%%%%%%%%

%\section*{Resumen ejecutivo}\addcontentsline{toc}{section}{Resumen ejecutivo} % Max 2 paginas, 5 puntos

%Ejemplo de sección sin numerar, indexada en el índice

%%%%%%%%%%%%%%%%%%%%%%%%%%%%%%%%%%%%%%%%%%%%%%%%%%%%%%%%%%%%%%%%%%%%%%%%%%%%%%%%%%%%%%%%%%%%%%%%%%
\newpage

\section{Introducción}%Máximo 2 paginas

Para poder utilizar palabras del glosario hay diferentes formatos:
\begin{itemize}
    \item \textbackslash gls\{\}: Ejemplo \gls{holding} palabra en singular con la letra inicial en minúscula.
    \item \textbackslash Gls\{\}: Ejemplo \Gls{holding} palabra en singular con la letra inicial en mayúscula.
    \item \textbackslash glspl\{\}: Ejemplo \glspl{holding} palabra en plural con la letra inicial en minúscula.
    \item \textbackslash Glspl\{\}: Ejemplo \Glspl{holding} palabra en plural con la letra inicial en mayúscula.
    
\end{itemize}

Para realizar las citas se utiliza el comando \textbackslash autocite y el nombre asignado en la bibliografía. \autocite{cencosud}

Para poner frases entre \q{comillas} cree el comando \textbackslash q\{\}

\noindent Para ignorar la sangría se realiza con \textbackslash noident al inicio del párrafo.

Ejemplo para insertar imágenes, gráficas, etc.

\begin{figure}[ht] 
    \centering
    \includegraphics[scale = 0.6]{IMAGENES/logo.png}
    \caption[Logo EII]{Logo EII. Fuente: eii.pucv.cl}

\end{figure}

Ejemplo para insertar tabla.
Recomiendo utilizar este sitio para generar tablas fácilmente.
\href{https://www.tablesgenerator.com/}{https://www.tablesgenerator.com/}

\begin{table}[H]
\small
\caption{Nemotécnicos, razones sociales e industrias de cada acción.}
\begin{tabular}{|l|l|l|}
\hline
\textbf{NEMOTÉCNICO} & \textbf{RAZÓN SOCIAL}                      & \textbf{INDUSTRIA} \\ \hline
SQM-B                & SOC QUIMICA   MINERA DE CHILE S.A. SERIE B & COMMODITIES        \\ \hline
CHILE                & BANCO DE CHILE                             & BANCA              \\ \hline
ENELAM               & ENEL AMERICAS   S.A.                       & UTILITIES          \\ \hline
\end{tabular}
\centering
\footnotesize Fuente: Rankia Chile, 2021
\end{table}


\newpage
\section{Antecedentes de la Empresa} %Máximo 5 paginas
En esta sección se describirán los principales antecedentes de la empresa donde el estudiante realizó su practica, tales como el proceso productivo de la empresa, sus principales productos, recursos y departamentos.
\subsection{Descripción del proceso productivo}

%\pagebreak

\subsection{Descripción de los principales productos}

Ejemplo de lista alfabética según formato escuela.

 \begin{enumerate}[label=\alph*.]
 \item Item a: ....
 \item Item b: ....
 \item Item c: ....
 \end{enumerate}
 
%\pagebreak

\subsection{Descripción de las principales materias primas o recursos}

Otro ejemplo de lista.

\begin{enumerate}[label=\alph*.]
\item \textbf{Item a}: ....
\item \textbf{Item b}: ....
\item \textbf{Item c}: ....
\end{enumerate}
 
%\pagebreak

\subsection{Descripción de los principales departamentos}

\newpage

\section{Antecedentes del departamento donde se realizó la práctica} % Máximo 3 paginas
En esta sección se describirán los detalles del departamento donde el estudiante realizó su practica, describiendo las principales actividades que se realizaron durante la estadía del estudiante, así como también los principales recursos utilizados en el departamento.

\subsection{Descripción de las principales actividades que se realizan}
Ejemplo de lista enumerada.

\begin{enumerate}
    \item Item 1 del proceso.
    \item Item 2 del proceso.
    \item Item 3 del proceso.
\end{enumerate}



\subsection{Descripción de los principales recursos}



%\pagebreak

\newpage

\section{Trabajo realizado}%Máximo 3 paginas
En esta sección se describirá el trabajo realizado por el estudiante durante su estadía en la empresa, así como también el ambiente en que el estudiante estuvo inmerso, también se realiza una breve descripción de la persona que acompañó laboralmente la estadía del estudiante.
\subsection{Descripción de las principales actividades realizadas}
\subsection{Descripción de la persona que acompañó laboralmente en la práctica}
\subsection{Descripción detallada del clima laboral y desempeño laboral del entorno de trabajo}

%\pagebreak

\newpage

\section{Conclusiones}%Máximo 3 paginas

La realización de su practica en XXXXX, trajo consigo una serie de nuevos logros y aprendizajes para el estudiante, los cuales serán descritos a continuación.

\subsection{Relacionadas al trabajo específico}
Respecto al trabajo especifico realizado, el estudiante logró ......

\subsection{Relacionadas a la experiencia personal}
Respecto a la experiencia personal se puede concluir que....

\subsection{Relacionadas con el cumplimiento de los objetivos}

Al termino de este periodo de practica, el estudiante ...



%%%%%%%%%%%%%%%%%%%%%%%%%%%%%%%%%%%%%%%%%%  Bibliografía %%%%%%%%%%%%%%%%%%%%%%%%%%%%%%%%%%%%% 

\newpage
\addcontentsline{toc}{section}{Bibliografía}
\printbibliography
 
 %%%%%%%%%%%%%%%%%%%%%%%%%%%%%%%%%%%%%%% Inicio Anexos %%%%%%%%%%%%%%%%%%%%%%%%%%%%%%%%%%%%%%%

\newpage
\hspace{0pt}
\vfill

\section*{\centering Anexos}\addcontentsline{toc}{section}{Anexos}
\vfill
\hspace{0pt}

%%%%%%%%%%%%%%%%%%%%%%%%%%%%%%%%%%%%%%%%%% Anexos %%%%%%%%%%%%%%%%%%%%%%%%%%%%%%%%%%%%%%%%%%

\section*{\centering Anexo 1 - Nombre del Anexo 1}\addcontentsline{toc}{section}{Anexo 1 - Nombre del Anexo 1}
\begin{table}[H]
\renewcommand\thetable{A1.1} 
\small
\caption{Nemotécnicos, razones sociales y la industria a la que pertenece cada acción. }
\begin{tabular}{|l|l|l|}
\hline
\textbf{NEMOTÉCNICO} & \textbf{RAZÓN SOCIAL}                      & \textbf{INDUSTRIA} \\ \hline
SQM-B                & SOC QUIMICA   MINERA DE CHILE S.A. SERIE B & COMMODITIES        \\ \hline
CHILE                & BANCO DE CHILE                             & BANCA              \\ \hline
ENELAM               & ENEL AMERICAS   S.A.                       & UTILITIES          \\ \hline
COPEC                & EMPRESAS COPEC   S.A.                      & UTILITIES          \\ \hline
BSANTANDER           & BANCO   SANTANDER-CHILE                    & BANCA              \\ \hline
CMPC                 & EMPRESAS CMPC   S.A.                       & COMMODITIES        \\ \hline
FALABELLA            & FALABELLA S.A.                             & RETAIL             \\ \hline
CENCOSUD             & CENCOSUD S.A.                              & RETAIL             \\ \hline
BCI                  & BANCO DE CREDITO   E INVERSIONES           & BANCA              \\ \hline
ENELCHILE            & ENEL CHILE S.A.                            & UTILITIES          \\ \hline
COLBUN               & COLBUN S.A.                                & UTILITIES          \\ \hline
PARAUCO              & PARQUE ARAUCO   S.A.                       & CONST. \& INMOB.   \\ \hline
CCU                  & COMPANIA   CERVECERIAS UNIDAS S.A.         & CONSUMO            \\ \hline
CAP                  & CAP S.A.                                   & COMMODITIES        \\ \hline
AGUAS-A              & AGUAS ANDINAS   S.A., SERIE A              & UTILITIES          \\ \hline
ENTEL                & EMP. NACIONAL DE   TELECOMUNICACIONES S.A. & COM \& TEC         \\ \hline
CONCHATORO           & VINA CONCHA Y   TORO S.A.                  & CONSUMO            \\ \hline
ANDINA-B             & EMBOTELLADORA   ANDINA S.A. SERIE B        & CONSUMO            \\ \hline
CENCOSHOPP           & CENCOSUD   SHOPPING S.A.                   & CONST. \& INMOB.   \\ \hline
VAPORES              & COMPANIA SUD   AMERICANA DE VAPORES S.A.   & INDUSTRIAL         \\ \hline
ECL                  & ENGIE ENERGIA   CHILE S.A.                 & UTILITIES          \\ \hline
ITAUCORP             & ITAU CORPBANCA                             & BANCA              \\ \hline
AESGENER             & AES GENER S.A.                             & UTILITIES          \\ \hline
MALLPLAZA            & PLAZA S.A.                                 & CONST. \& INMOB.   \\ \hline
IAM                  & INVERSIONES   AGUAS METROPOLITANAS S.A.    & UTILITIES          \\ \hline
RIPLEY               & RIPLEY CORP S.A.                           & RETAIL             \\ \hline
SMU                  & SMU S.A.                                   & RETAIL             \\ \hline
SONDA                & SONDA S.A.                                 & COM \& TEC         \\ \hline
ILC                  & INVERSIONES LA   CONSTRUCCION S.A.         & CONST. \& INMOB.   \\ \hline
SECURITY             & GRUPO SECURITY   S.A.                      & BANCA              \\ \hline
\end{tabular}
\centering
\footnotesize Fuente: Rankia Chile, 2021
\end{table}
\newpage
%%%%%%%%%%%%%%%%%%%%%%%%

\section*{\centering Anexo 2 - Nombre del Anexo 2}\addcontentsline{toc}{section}{Anexo 2 - Nombre del Anexo 2}
\begin{table}[H]
\renewcommand\thetable{A2.1} 
\small
\caption{Precios de apertura y cierre de IPSA, SQM-B Y AGUAS-A.}
\begin{tabular}{|c|c|c|c|c|c|}
\hline
\multicolumn{2}{|c|}{IPSA}                       & \multicolumn{2}{c|}{SQM-B}                           & \multicolumn{2}{c|}{AGUAS-A}                          \\ \hline
\textit{\textbf{Open}} & \textit{\textbf{Close}} & \textit{\textbf{Open}} & \textit{\textbf{Close}} & \textit{\textbf{Open}} & \textit{\textbf{Close}} \\ \hline
4253,9126 & 4252,043  & 37000 & 36250 & 145,01 & 154,3  \\ \hline
4252,043  & 4211,6963 & 36250 & 36475 & 153,99 & 152,8  \\ \hline
4211,6963 & 4207,026  & 36475 & 37402 & 152,8  & 150    \\ \hline
4207,026  & 4211,4756 & 37601 & 37498 & 150    & 153,9  \\ \hline
4211,4756 & 4240,5444 & 37498 & 38510 & 155,33 & 153,78 \\ \hline
4240,5444 & 4282,479  & 38531 & 39489 & 153,78 & 153,31 \\ \hline
4282,479  & 4243,715  & 39491 & 39900 & 149    & 151,41 \\ \hline
4243,715  & 4345,423  & 39900 & 41800 & 151,41 & 151,97 \\ \hline
4345,423  & 4367,996  & 42382 & 40896 & 151,97 & 152,14 \\ \hline
4367,996  & 4374,7544 & 40705 & 41102 & 152,14 & 153    \\ \hline
4374,7544 & 4401,936  & 40897 & 42001 & 153    & 152,99 \\ \hline
4401,936  & 4329,4146 & 41300 & 40375 & 152,99 & 151,08 \\ \hline
4329,4146 & 4320,577  & 39999 & 38850 & 151,08 & 152,99 \\ \hline
4320,577  & 4341,585  & 39000 & 38490 & 149    & 153,99 \\ \hline
4341,585  & 4310,5103 & 37500 & 37700 & 153,99 & 155,5  \\ \hline
4310,5103 & 4347,4854 & 37901 & 38949 & 155,5  & 159    \\ \hline
4347,4854 & 4392,7485 & 39552 & 39803 & 159,8  & 166,8  \\ \hline
4392,7485 & 4369,432  & 40202 & 39550 & 167,28 & 167    \\ \hline
4369,432  & 4437,935  & 39549 & 40885 & 167    & 166,15 \\ \hline
4437,935  & 4443,8057 & 40880 & 40105 & 166,15 & 174,45 \\ \hline
4443,8057 & 4459,2173 & 40100 & 40240 & 174,45 & 172    \\ \hline
4459,2173 & 4506,5337 & 40650 & 40750 & 172    & 173,03 \\ \hline
4506,5337 & 4491,5723 & 40750 & 40481 & 173,03 & 173,13 \\ \hline
4491,5723 & 4450,955  & 39900 & 40351 & 173    & 172,66 \\ \hline
4450,955  & 4437,4404 & 39999 & 41818 & 172,66 & 171,79 \\ \hline
4437,4404 & 4463,5527 & 41818 & 41790 & 171,79 & 174,5  \\ \hline
4463,5527 & 4456,2036 & 41780 & 41350 & 174,5  & 172,02 \\ \hline
4456,2036 & 4432,984  & 42000 & 41395 & 172,02 & 171    \\ \hline
4432,984  & 4378,1436 & 41500 & 41324 & 171    & 167,25 \\ \hline
4378,1436 & 4407,236  & 41324 & 42650 & 167,25 & 166    \\ \hline
4407,236  & 4463,9995 & 43001 & 44301 & 166    & 166,8  \\ \hline
\end{tabular}
\centering
\end{table}
\newpage
%%%%%%%%%%%%%%%%%%%%%%%%

\section*{\centering Anexo 3 - Nombre del Anexo 3}\addcontentsline{toc}{section}{Anexo 3 - Nombre del Anexo 3}
 
\begin{table}[H]
\renewcommand\thetable{A3.1} 
\caption{Precios de apertura y cierre de IPSA, SQM-B Y AGUAS-A y sus respectivas variaciones porcentuales.}
\begin{tabular}{|c|c|c|c|c|c|c|c|c|}
\hline
\multicolumn{3}{|c|}{IPSA} & \multicolumn{3}{c|}{SQM-B} & \multicolumn{3}{c|}{AGUAS-A} \\ \hline
\textit{\textbf{Open}} & \textit{\textbf{Close}} & \textit{\textbf{$\Delta \%$}} & \textit{\textbf{Open}} & \textit{\textbf{Close}} & \textit{\textbf{$\Delta \%$}} & \textit{\textbf{Open}} & \textit{\textbf{Close}} & \textit{\textbf{$\Delta \%$}} \\ \hline
4253,9126 & 4252,043 & -0,044 & 37000 & 36250 & -2,027 & 145,01 & 154,3 & 6,406 \\ \hline
4252,043 & 4211,6963 & -0,949 & 36250 & 36475 & 0,621 & 153,99 & 152,8 & -0,773 \\ \hline
4211,6963 & 4207,026 & -0,111 & 36475 & 37402 & 2,541 & 152,8 & 150 & -1,832 \\ \hline
4207,026 & 4211,4756 & 0,106 & 37601 & 37498 & -0,274 & 150 & 153,9 & 2,600 \\ \hline
4211,4756 & 4240,5444 & 0,690 & 37498 & 38510 & 2,699 & 155,33 & 153,78 & -0,998 \\ \hline
4240,5444 & 4282,479 & 0,989 & 38531 & 39489 & 2,486 & 153,78 & 153,31 & -0,306 \\ \hline
4282,479 & 4243,715 & -0,905 & 39491 & 39900 & 1,036 & 149 & 151,41 & 1,617 \\ \hline
4243,715 & 4345,423 & 2,397 & 39900 & 41800 & 4,762 & 151,41 & 151,97 & 0,370 \\ \hline
4345,423 & 4367,996 & 0,519 & 42382 & 40896 & -3,506 & 151,97 & 152,14 & 0,112 \\ \hline
4367,996 & 4374,7544 & 0,155 & 40705 & 41102 & 0,975 & 152,14 & 153 & 0,565 \\ \hline
4374,7544 & 4401,936 & 0,621 & 40897 & 42001 & 2,699 & 153 & 152,99 & -0,007 \\ \hline
4401,936 & 4329,4146 & -1,647 & 41300 & 40375 & -2,240 & 152,99 & 151,08 & -1,248 \\ \hline
4329,4146 & 4320,577 & -0,204 & 39999 & 38850 & -2,873 & 151,08 & 152,99 & 1,264 \\ \hline
4320,577 & 4341,585 & 0,486 & 39000 & 38490 & -1,308 & 149 & 153,99 & 3,349 \\ \hline
4341,585 & 4310,5103 & -0,716 & 37500 & 37700 & 0,533 & 153,99 & 155,5 & 0,981 \\ \hline
4310,5103 & 4347,4854 & 0,858 & 37901 & 38949 & 2,765 & 155,5 & 159 & 2,251 \\ \hline
4347,4854 & 4392,7485 & 1,041 & 39552 & 39803 & 0,635 & 159,8 & 166,8 & 4,380 \\ \hline
4392,7485 & 4369,432 & -0,531 & 40202 & 39550 & -1,622 & 167,28 & 167 & -0,167 \\ \hline
4369,432 & 4437,935 & 1,568 & 39549 & 40885 & 3,378 & 167 & 166,15 & -0,509 \\ \hline
4437,935 & 4443,8057 & 0,132 & 40880 & 40105 & -1,896 & 166,15 & 174,45 & 4,995 \\ \hline
4443,8057 & 4459,2173 & 0,347 & 40100 & 40240 & 0,349 & 174,45 & 172 & -1,404 \\ \hline
4459,2173 & 4506,5337 & 1,061 & 40650 & 40750 & 0,246 & 172 & 173,03 & 0,599 \\ \hline
4506,5337 & 4491,5723 & -0,332 & 40750 & 40481 & -0,660 & 173,03 & 173,13 & 0,058 \\ \hline
4491,5723 & 4450,955 & -0,904 & 39900 & 40351 & 1,130 & 173 & 172,66 & -0,197 \\ \hline
4450,955 & 4437,4404 & -0,304 & 39999 & 41818 & 4,548 & 172,66 & 171,79 & -0,504 \\ \hline
4437,4404 & 4463,5527 & 0,588 & 41818 & 41790 & -0,067 & 171,79 & 174,5 & 1,578 \\ \hline
4463,5527 & 4456,2036 & -0,165 & 41780 & 41350 & -1,029 & 174,5 & 172,02 & -1,421 \\ \hline
4456,2036 & 4432,984 & -0,521 & 42000 & 41395 & -1,440 & 172,02 & 171 & -0,593 \\ \hline
4432,984 & 4378,1436 & -1,237 & 41500 & 41324 & -0,424 & 171 & 167,25 & -2,193 \\ \hline
4378,1436 & 4407,236 & 0,664 & 41324 & 42650 & 3,209 & 167,25 & 166 & -0,747 \\ \hline
4407,236 & 4463,9995 & 1,288 & 43001 & 44301 & 3,023 & 166 & 166,8 & 0,482 \\ \hline
\end{tabular}
\centering
\end{table}

\newpage
%%%%%%%%%%%%%%%%%%%%%%%%

\section*{\centering Anexo 4 - Nombre del Anexo 4}\addcontentsline{toc}{section}{Anexo 4 - Nombre del Anexo 4}

\begin{figure}[ht]
\renewcommand\thefigure{A4.1} 
    \centering
    \includegraphics[scale = 0.7]{IMAGENES/logo.png}
    \caption{Logo EII.}
    \label{CodigoDescribe}
\end{figure}

%%%%%%%%%%%%%%%%%%%%%%%%


\begin{figure}[ht]
\renewcommand\thefigure{A4.2} 
    \centering
    \includegraphics[scale = 0.7]{IMAGENES/logo.png}
    \caption{Logo EII.}
    \label{CodigoCovVar}
\end{figure}

 
%%%%%%%%%%%%%%%%%%%%%%%%



\begin{figure}[ht]
\renewcommand\thefigure{A4.3} 
    \centering
    \includegraphics[scale = 0.9]{IMAGENES/logo.png}
    \caption{Logo EII.}
    \label{CodigoTests}
\end{figure}
 

\end{document}


%Conversor de formulas a LaTeX
% https://www.codecogs.com/latex/eqneditor.php 


\end{document}